\booksection{제1장 부록: HarmonyOS 운영체제 특집}

\chapterintro{HarmonyOS는 화웨이가 자체 개발한 운영체제로, 2019년 출시 이후 여러 메이저 버전을 거쳐 왔습니다. 이것은 단순한 “또 다른 안드로이드 대체품”이 아니라, 독특한 설계 철학을 가진 새로운 종입니다. 이 부상 중인 운영체제를 함께 깊이 알아봅시다.}

\booksubsection{HarmonyOS란 무엇인가? 왜 그것은 다른가?}

HarmonyOS의 핵심 철학은 “사물 인터넷”입니다. 전통적인 운영체제가 주로 단일 기기에 서비스를 제공하는 것과 달리, HarmonyOS는 설계 초기부터 다중 기기 협동을 고려했습니다.

\textbf{분산 아키텍처}는 HarmonyOS의 영혼입니다. 간단히 말해, 여러 기기를 하나의 “슈퍼 터미널”로 가상화할 수 있습니다. 스마트폰, 태블릿, 컴퓨터, 시계, TV는 더 이상 고립된 개체가 아니라, 전체의 “장기”로서 서로의 능력을 서로 호출할 수 있습니다.

예를 들어: 당신은 스마트폰으로 영상을 보다가 대화면에서 보고 싶어지면, 전통 방식은 화면 미러링; 반면 홍몽에서는 영상을 TV로 “이동”시켜 계속 볼 수 있고, 스마트폰은 리모컨이 됩니다. 이것은 간단한 화면 미러링이 아니라, 앱 자체가 기기 간 실행되는 것입니다.

\booksubsection{홍몽 4.3: AOSP 기반의 성숙 버전}

홍몽 4.3은 현재 화웨이 스마트폰에 광범위하게 배치된 버전입니다. 설명해야 할 것은, 홍몽 4.3은 안드로이드 오픈소스 프로젝트(AOSP)의 코드 프레임워크와 호환되므로, 안드로이드 앱을 완전히 호환할 수 있습니다 - 당신이 화웨이 앱 마켓에서 다운로드한 대부분의 App은 본질적으로 여전히 안드로이드 앱이며, 홍몽의 호환 층을 통해 실행됩니다. 단 홍몽 4.3은 간단한 겉껍데기 안드로이드가 아니며, AOSP 기반 위에대량 화웨이 자체 개발 분산 능력, 방주 엔진 등 기술을 추가했습니다.

\subsubsection*{핵심 특성:}

\begin{itemize}
  \item \textbf{안드로이드 앱 호환}: AOSP 기반으로 개발되었기 때문에, HarmonyOS 4.3은 거의 모든 안드로이드 앱을 정상적으로 실행할 수 있으며, 사용자는 거의 차이를 느끼지 못함.
  \item \textbf{방주 엔진}: 화웨이 자체 개발 컴파일러와 런타임으로, 앱 실행을 더 매끄럽고 더 절전하게 만듦. 초기 버전 대비 유창도 향상이 뚜렷함.
  \item \textbf{슈퍼 터미널}: 제어 센터에서 주변의 다른 홍몽 기기를 직관적으로 볼 수 있으며, 드래그만 하면 연결. 스마트폰이 컴퓨터의 키보드·마우스를 호출하거나, 태블릿이 스마트폰의 카메라를 호출하는 것이 모두 편리함.
  \item \textbf{서비스 카드}: iOS의 위젯과 유사하지만 더 유연함. App 아이콘을 위로 슬라이드하여 카드를 펼치고, 핵심 정보를 표시하며, App을 열 필요 없음.
  \item \textbf{순수 모드}: 공식 앱 마켓에서만 앱 설치를 허용, 악성 소프트웨어 위험 감소.
  \item \textbf{프라이버시 센터}: 앱 권한을 집중 관리하고, 앱 행위 기록을 조회.
\end{itemize}

\subsubsection*{사용자 인터페이스 특징:}

\begin{itemize}
  \item \textbf{매거진 잠금 화면}: 잠금 화면에 정교한 이미지와 기사를 표시하며, 쓸어 넘기면 즉시 데스크톱으로 진입.
  \item \textbf{만능 카드}: 데스크톱 아이콘을 카드로 펼칠 수 있어, 더 많은 정보 표시.
  \item \textbf{슈퍼 중계소}: 임시 저장 영역으로, 문자, 사진, 파일을 임시 저장한 후 임의의 App으로 드래그하여 사용. 앱 간 공유가 매우 매끄러워짐.
  \item \textbf{스마트 검색}: 전역 검색으로, 앱, 파일, 설정, 웹페이지 검색 가능.
\end{itemize}

\booksubsection{HarmonyOS NEXT (HarmonyOS 6): 화웨이 자체 운영체제의 진정한 시작}

HarmonyOS NEXT (HarmonyOS NEXT 5.0)는 HarmonyOS 4.3과 근본적인 차이가 있습니다 - 화웨이 자체 개발 HarmonyOS 마이크로 커널을 채택하여, 더 이상 Linux 커널과 AOSP 코드에 의존하지 않으며,\textbf{더 이상 안드로이드 앱을 직접 지원하지 않습니다}. 이는 HarmonyOS가 “안드로이드 호환 과도기 형태”에서 “완전히 독립적인 운영체제”로 공식적으로 나아가는 중요한 단계입니다.

\subsubsection*{중대한 변혁:}

\begin{itemize}
  \item \textbf{완전 자주 커널}: 홍몽 NEXT는 화웨이 자체 개발 커널 사용, 더 이상 어떠한 AOSP 코드도 포함하지 않음. 시스템이 더 간결하고, 더 안전하고, 더 효율적.
  \item \textbf{안드로이드 앱 더 이상 지원 안 함}: 이것이 가장 큰 변화. 홍몽 NEXT는 원생 홍몽 앱만 실행할 수 있으며, APK 파일 설치 및 실행 불가. 업그레이드 전 당신이 자주 쓰는 App이 이미 홍몽 원생 버전이 있는지 확인하세요.
  \item \textbf{원생 앱 생태}: 화웨이는 개발자가 앱을 홍몽 원생 플랫폼으로 이전하도록 전력 추진 중. 현재 주류 App인 위챗, 알리페이, 타오바오 등은 이미 홍몽 원생 버전을 출시했으나, 일부 마이너 앱은 아직 적응되지 않았을 수 있음.
  \item \textbf{성방 보안 아키텍처}: 새로운 보안 체계, 하드웨어에서 소프트웨어까지 전 링크 보호. 더 세밀한 권한 제어 지원, 예: "퍼지 위치" - 앱은 당신이 대략 어디에 있는지만 알 수 있고, 정확한 위치는 알 수 없음.
  \item \textbf{AI 심층 통합}: 시스템 레벨 AI 능력, 단지 어떤 앱이 AI를 가지는 것이 아님. 예: AI 자막은 영상을 실시간 번역, AI 컷아웃은 피사체를 지능적으로 식별.
  \item \textbf{기기 간 매끄러운 전환}: 더 한 단계 발전한 분산 능력. 앱 상태가 기기 간에 매끄럽게 전환 가능, 스마트폰에서 편집한 문서가 컴퓨터 앞으로 이동하면 자동으로 컴퓨터에서 열림.
  \item \textbf{메타 서비스}: 설치 없이 사용 가능한 경량 서비스. 예: QR 코드 스캔으로 즉시 어떤 기능 열기, 먼저 앱 다운로드 불필요.
\end{itemize}

\begin{tcbwarn}{HarmonyOS NEXT 업그레이드 전 필독}
HarmonyOS NEXT로 업그레이드하면, 스마트폰은 더 이상 안드로이드 앱을 실행할 수 없습니다. 업그레이드 전에 반드시 확인하세요: (1) 자주 사용하는 앱이 이미 HarmonyOS 네이티브 버전이 있는지; (2) 업무나 생활에 반드시 안드로이드 앱에 의존해야 하는 시나리오가 없는지. 업그레이드 후 HarmonyOS 4.3으로 돌아갈 수 없음.

\end{tcbwarn}

\subsubsection*{작업 논리 변화:}

\begin{itemize}
  \item \textbf{제스처 내비게이션}: 좌우 가장자리 스와이프로 뒤로, 하단 상단 스와이프로 데스크톱, 상단 스와이프 멈춤으로 멀티태스크 진입.
  \item \textbf{스마트 멀티 윈도우}: 분할 화면, 부유 창, 평행 시야 등 다양한 멀티태스킹 형태 지원.
  \item \textbf{창련 통화}: 화웨이의 영상 통화 서비스, 화면 공유, 뷰티, 배경 흐림 지원.
\end{itemize}

\booksubsection{HarmonyOS 생태계: 기기와 앱}

홍몽은 스마트폰 시스템만이 아니라, 다양한 기기를 커버하는 생태 시스템입니다.

\subsubsection*{지원 기기 유형:}

\begin{itemize}
  \item \textbf{스마트폰과 태블릿}: 핵심 기기, 기능이 가장 완전.
  \item \textbf{스마트 워치와 밴드}: 독립 실행 또는 스마트폰과 협동.
  \item \textbf{스마트 스크린 (TV)}: 가정용 오락 센터.
  \item \textbf{스마트 스피커}: 음성 상호작용 진입점.
  \item \textbf{자동차}: HarmonyOS 콕핏, 지능형 운전석 경험.
  \item \textbf{스마트 홈}: 냉장고, 에어컨, 세탁기 등.
\end{itemize}

\subsubsection*{앱 생태계 현황:}

HarmonyOS 앱 마켓에는 이미 수백만 개의 앱이 있으며, 주류 앱은 기본적으로 모두 있습니다. 화웨이도 개발자들이 네이티브 HarmonyOS 앱을 개발하도록 적극적으로 추진 중. 2026년 중반까지, 위챗, 알리페이, 타오바오, 도우인, 메이퇀, 징둥 등 주요 앱은 모두 HarmonyOS 네이티브 버전을 출시하여 일상 사용에 기본적으로 무방. 일부 마이너 앱 및 업계 앱은 여전히 적응 중일 수 있으며, 업그레이드 전에 자주 사용하는 앱이 이미 지원되는지 확인하는 것을 권장.5. HarmonyOS 실용 팁\textbf{슈퍼 터미널}: 화면 우측 상단에서 아래로 쓸어밀어 제어 센터를 열고, “슈퍼 터미널”을 클릭하면 주변의 다른 기기를 볼 수 있으며, 드래그하여 연결.\textbf{슈퍼 중계소}: 텍스트/이미지/파일을 길게 누르고 “중계소로 복사”를 선택한 다음, 다른 앱에서 길게 누르고 붙여넣기 또는 드래그하여 꺼내기.\textbf{서비스 카드}: 앱 아이콘을 위로 쓸어밀어 카드를 펼치고, 카드를 길게 누르면 크기 조정 또는 데스크톱에 추가 가능.\textbf{스마트 멀티 윈도우}: 화면 좌우 가장자리에서 안쪽으로 쓸어밀고 멈추어 사이드바를 열고, 앱을 클릭하여 부유 창으로 열기.\textbf{프라이버시 보호}: 설정 > 프라이버시, 앱 권한 사용 기록 조회, 퍼지 위치 등 설정 가능.주의 사항HarmonyOS 시스템은 계속 업데이트되고 발전 중이며, 버전에 따라 약간의 차이가 있을 수 있음. 본문의 설명은 HarmonyOS 4.3과 6의 주요 특성을 기반으로 하며, 구체적인 기능은 스마트폰의 실제 버전을 기준으로 해야 함. 최신 기능 및 보안 패치를 받기 위해 시스템 업데이트를 유지하는 것을 권장.지금 해보기: HarmonyOS 특징 기능 체험제어 센터를 열고, “슈퍼 터미널”을 찾아 주변에 다른 HarmonyOS 기기가 있는지 확인.“슈퍼 중계소” 시도: 텍스트 한 조각을 중계소에 복사한 다음, 메모장에서 드래그하여 꺼내기.자주 사용하는 앱 아이콘을 몇 개 위로 쓸어밀어 서비스 카드에 어떤 내용이 표시되는지 확인.설정 > 프라이버시로 이동하여 최근 어떤 앱이 민감 권한을 사용했는지 확인.
